\section{Introduction}\label{sec:intro}

% G. General
With the rapid growth of open government data and urban sensing, municipal services increasingly publish operational data in real time. In many Taiwanese cities, household waste is not left at the curb but handed directly to a collection truck that stops briefly at designated points according to a published schedule. Taichung City publishes the positions of its garbage and recycling trucks as open data, with records containing route, vehicle, timestamp, location, coordinates, and speed fields refreshed approximately every ten minutes~\cite{taichung_fleet_api}. Such data make it possible to see \emph{where} a truck is, and they have enabled map-based tracking services for residents.

% H. However
However, residents face a different question: \emph{given where I am and how fast I walk, should I leave now, and for which stop?} This question is becoming more consequential as the population ages. Taiwan crossed the threshold of a super-aged society, in which more than 20\% of the population is aged 65 or above, in 2025~\cite{ndc_population}, and walking speed varies substantially with age~\cite{knoblauch1996}. For an older resident, leaving too early means waiting at the curb, while leaving too late means missing the truck and storing waste for another day. Answering the question requires combining three uncertain quantities: the truck's arrival time at a stop, the resident's walking time to that stop, and the freshness of the data on which both estimates depend. With location updates only every ten minutes, none of these can be made exact, and an optimistic ``you can make it'' is more harmful than no advice at all.

% I. In literature
In the literature, solutions related to this problem have evolved from simple to complex. The earliest are \emph{static schedules}, which give a nominal collection time per stop. \emph{Vehicle tracking displays} add the latest GPS position. \emph{Arrival-time prediction models}, developed mainly for public buses, estimate arrivals from historical averages, regression, or neural networks~\cite{chien2002,jeong2004}. \emph{Pedestrian routing} services estimate walking time over road networks~\cite{postgis}. Most recently, \emph{tool-using large language model (LLM) agents} can interpret natural-language requests and call external services~\cite{yao2023react,schick2023toolformer}.

% J. Judgement
Each family has clear strengths and weaknesses. Static schedules and tracking displays are simple and transparent, but they do not intersect truck arrival with walking time. Arrival-time models improve point accuracy but assume dense observations and do not decide whether a resident should depart. Pedestrian routing ignores collection-route logic and dwell behavior at stops. LLM agents offer natural interaction, but without grounding they may produce plausible yet fabricated times and distances~\cite{ji2023hallucination}. None of these approaches decides \emph{when it is safe to advise and when it should abstain}, a question formalized by the reject option~\cite{chow1970} and selective prediction~\cite{geifman2017} but not yet applied to curbside collection.

% K. Keypoint
The key idea of this work is a paradigm shift from \emph{predicting arrivals more accurately} to \emph{deciding when a departure recommendation is justified}. We present \emph{GarbageGo}, a domain-specific, tool-grounded agent for departure decisions in curbside waste collection in Taichung City. Its contributions are as follows:
\begin{itemize}
\item \textbf{An advise-or-defer decision formulation.} Departure feasibility is defined by a time margin (slack) between predicted truck arrival, estimated walking time, and a safety buffer that widens with data latency, prediction error, and positioning uncertainty, with explicit conditions under which the system must defer. We show that a calibrated buffer bounds the false-feasibility rate.
\item \textbf{A tool-grounded agent architecture.} Domain knowledge, including schedules, pedestrian accessibility, and collection-operation characteristics, is encapsulated as whitelisted tools with JSON Schema-validated inputs and outputs. The LLM is restricted to intent recognition and explanation, and every number it reports must originate from a tool result.
\item \textbf{A safety-oriented evaluation protocol.} A rolling-origin backtest compares arrival-time predictors, treats the false-feasibility rate as the primary safety metric alongside recommendation coverage and fallback success, and includes at least thirty anomaly scenarios.
\end{itemize}
This paper reports the formulation, system design, and evaluation protocol. Empirical results will be reported after the backtest is completed; no experimental values are claimed here.

% L. List the organization
The remainder of this paper is organized as follows. Section~\ref{sec:related} reviews and compares related methods. Section~\ref{sec:method} presents the proposed scheme. Section~\ref{sec:experiments} describes the simulation and evaluation design, and Section~\ref{sec:conclusion} concludes the paper.
